
\ARMTable{}


% What is the exchangeable unit?
As discussed in \secref{arm_experiments}, for the ARM experiments we needed to
make decisions about what was an exchangeable unit. For models with no random
effects, we define an exchangeable unit to be a tuple containing a particular
response and a regressor.  For random effects models, we had to choose an
exchangeable unit. When the random effect had many distinct levels, we
considered the groups of observations within a level to be a single exchangeable
unit.  For models whose random effects had few distinct levels, we took the
exchangeable unit to either be a more fine-grained partition of the data, or, in
some cases treated the observations as fully independent.

Although the posterior covariance is not of primary
interest, we compute it to investigate the possibility of differences between
the frequentist and posterior covariance estimates.
%
Given that most of the models in the ARM textbook are carefully selected by
highly experienced statisticians, one might expect there to be few instances of
misspecification.  Since the models are additionally mostly low-dimensional, we
might expect for the posterior covariance to typically provide a good
approximation to the frequentist covariance by the Bernstein von-Mises theorem.

To investigate this intuition, in \figref{normerr_post_graph} we reproduce
\figref{normerr_graph} but comparing the bootdstrap and Bayesian posterior
covariance.  The metric plotted is the posterior version of
$\normdiffij_{ij}$ in \eqref{arm_metrics}:
$$
\normdiffbayes_{ij} :=
  \frac{\gcovbayeshat_{ij} - \gcovboothat_{ij}}
  {\left| \gcovboothat_{ij} \right| + \seboot_{ij}}.
$$

\ARMGraphPostDiff{}

Comparing with \figref{normerr_graph}, we see that the above intuition is mostly
correct, and that the frequentist and posterior covariances are largely similar.
The notable exception is for log scale parameters in models with relatively
large datasets, for which the IJ provides a markedly better approximation to the
bootstrap than does the posterior covariance.  This suggests the possibility of
using the difference between IJ and posterior covariances as a diagnostic for
model misspecification, as is sometimes done in the frequentist literature
\citep{white:1982:maximum}.  We leave this possibility as an interesting avenue
for future work.